
\section{Avaliação Experimental}
\label{sec:avaliacao}



%A fim de avaliar a proposta deste trabalho, conduziu-se os experimentos descritos a seguir. A Subseção~\ref{subsec:bases}, descreve como os 


%%%%%%%%%%
% As 2 principais dificuldades na avaliação
%%%%%%%%%%
De acordo com \cite{Pevzner2002} há duas principais dificuldades na avaliação de segmentadores automáticos. A primeira é conseguir uma referência, já que juízes humanos costumam não concordar entre si, sobre onde os limites estão e outras abordagens podem não se aplicar ao contexto. A segunda é que tipos diferentes de erros devem ter pesos diferentes de acordo com a aplicação. Há casos onde certa imprecisão é tolerável e outras, como a segmentação de notícias, onde a precisão é mais importante.


%%%%%%%%%%
% Definição do que é um bom algoritmo de segmentação
%%%%%%%%%%
Para fins de avaliação desse trabalho, um bom método de segmentação é aquele cujo resultado melhor se aproxima do ideal, sem a obrigatoriedade de estar perfeitamente alinhado com tal. Ou seja, visto o contexto das atas de reunião, e a subjetividade da tarefa, não é necessário que os limites entre os segmentos (real e hipótese) sejam idênticos, mas que se assemelhem em localização e quantidade.


%%%%%%%%%%
% Tratamento das atas
%%%%%%%%%%


As próximas subseções mostram o conjunto de atas e a segmentação usada como referência em seguida são apresentadas as métricas de avaliação utilizadas neste trabalho e os testes realizados para avaliar os métodos.



\subsection{Bases}
	\label{subsec:bases}

	
%%%%%%%%%%
% Como foram obtidas (software e especialistas)
%%%%%%%%%%
A fim de obter um conjunto de documentos segmentados que possam servir como referência na avaliação, foram coletadas as atas de reuniões do Conselhos de Pós Graduação, Conselho de Cursos e Conselho de Departamento de da UFSCar-Sorocaba. Os documentos foram oferecidos à profissionais que participam de reuniões desse departamento para segmentá-las. Para isso, utilizou-se um \textit{software} que permitiu aos voluntários visualizar um documento, e indicar livremente as divisões entre segmentos com um assunto relativamente independente. Ao final, o software coletou os dados da segmentação de doze das os quais serviram como referência para a avaliação dos algoritmos.

Os arquivos gerados foram tratados para que os segmentos sempre terminem em uma sentença reconhecida pelo algoritmo, uma vez que as sentenças são a unidade mínima de informação nesse trabalho.
  

\subsection{Configuração experimental}
	\label{subsec:configuracaoexperimental}

%%%%%%%%%%
% Parâmetros
%%%%%%%%%%
As implementações dos algoritmos permitem ao usuário a configuração de seus parâmetros. 
%
O \textit{TextTiling} permite ajustarmos dois parâmetros, sendo, o tamanho da janela (distância entre a primeira e a última sentença) para o qual atribuiu-se os valores 20, 40 e 60. Para o segundo parâmetro, o passo (distância que a janela desliza), atribuiu-se os valores 3, 6, 9 e 12. Gerando ao final 20 modelos.
%

O \textit{C99} permite ajustarmos três parâmetros, sendo, a quantidade segmentos desejados, o qual é calculado como uma proporção dos candidatos a limite. Para isso atribuiu-se as proporções {0,2; 0,4; 0,6; 0,8; 1,0}. O segundo parâmetro, o tamanho do quadro utilizado para gerar a matriz de ranking, atribuiu-se os valores 9 e 11, sendo 11 o valor padrão da apresentado pelo autor. Permite ainda, definirmos se as sentenças serão representados por vetores contendo a frequência ou o peso de cada termo, onde ambas as representações foram utilizadas.
%
 Considerando todos os parâmetros, foram gerados 20 modelos para o algoritmo C99.% By Rafael 


\subsection{Critérios de avaliação}


Os critérios vão aqui 



\subsection{Resultados}
	\label{subsec:resultados}
	
	
  \begin{figure}[!h]
	\centering
	\includegraphics[width=0.2\textwidth]{construcao.png}
%	\caption{Exemplo de pré-processamento.}
	\\Em Construção
	\label{fig:emconstrucao}
  \end{figure}

	

% seis atas de reunião foram coletadas junto ao departamento de pós-graduação da UFSCar-Sorocaba\footnote{\urlatas}. Os documentos foram oferecidos à profissionais que participam de reuniões desse departamento e por meio de um \textit{software} segmentaram o texto das atas conforme o julgamento de cada um. Os segmentos gerados manualmente foram comparados à segmentação automática conforme os critérios descritos a seguir.


%%%%%%%%%%
% Objetivos
%%%%%%%%%%
%A fim de encontrar o melhor método que divida uma ata em segmentos coerentes, realizou-se experimentos com o \textit{TextTiling} e \textit{C99} a fim de encontrar os melhores parâmetros para esses documentos.

%%%%%%%%%%
% Teste de Fiedman e CD
% 1ª Etapa
%%%%%%%%%%
%Em seguida aplicou-se o teste de Friedman para saber se há diferenças significativas entre a eficácia dos modelos. O pós-teste de Nemenyi foi aplicado para descobrir quais diferenças são significativas. 
%
%Exite diferença quando seus \textit{rankings} médios diferirem em um valor mínimo, chamado de diferença critica (CD). 
%

%%%%%%%%%%
% 
%%%%%%%%%%
%Com isso foi possível, pela análise do diagrama de diferença crítica, verificar qual é o melhor modelo para cada medida
% e quão significativamente 
%em relação aos demais. 

%%%%%%%%%%
% Dados obtidos com o C99
%%%%%%%%%%
%A Tabela~\ref{tab:mediasC99} mostra os dados obtidos com o \textit{C99}, onde \texttt{S} é a proporção de segmentos em relação a quantidade de candidatos, \texttt{M} é o tamanho do quadro utilizado para criar a matriz de \textit{ranking} e \texttt{W} indica se os segmentos são representados por vetores contendo a frequência ou um peso das palavras. 



%\begin{table}[!h]
%	\centering
%
%	\begin{tabular}{|c|c|c|c|c|}
%	
%		\hline
%		Medida & \texttt{S} & \texttt{M} & \texttt{W} & \textbf{Média}\\		
%		\hline
%
%		Acuracy		& 40	& 11 & Sim & 0.6199	\\ \hline	
%		F1			& 60	& 9	 & Sim & 0.6167	\\ \hline	
%		Precision	& 40	& 11 & Sim & 0.7106	\\ \hline			
%		Recall		& 100	& 9	 & Não & 0.8516	\\ \hline		
%		Pk			& 40	& 11 & Sim & 0.1163	\\ \hline	
%		Windiff		& 40	& 11 & Sim & 0.3800	\\ \hline		
%
%		
%	\end{tabular}
%	
%	\caption{Médias das medidas obtidas com \textit{C99}.}
%	\label{tab:mediasC99}
%\end{table}


%A tabelas~\ref{tab:mediasTextTiling} mostra os dados obtidos com o \textit{TextTiling}, onde \texttt{J} é o tamanho da janela e \texttt{P} é o passo.

%\begin{table}[!h]
%	\centering
%
%	\begin{tabular}{|c|c|c|c|}
%	
%		\hline
%		Medida & \texttt{J} & \texttt{P} & \textbf{Média}\\		
%		\hline
%
%		Acuracy		& 50 & 9 	& 0.5510 \\ \hline	
%		F1			& 50 & 3 	& 0.5898 \\ \hline	
%		Precision	& 60 & 12 	& 0.5746 \\ \hline			
%		Recall		& 50 & 3 	& 0.7717 \\ \hline		
%		Pk			& 30 & 9 	& 0.1572 \\ \hline	
%		Windiff		& 50 & 9 	& 0.4489 \\ \hline		
%
%		
%	\end{tabular}
%	
%	\caption{Médias das medidas obtidas com o \textit{TextTiling}.}
%	\label{tab:mediasTextTiling}
%\end{table}


%Uma vez sabendo quais valores de parâmetros melhor configuram um algoritmo para uma medida, resta então saber qual dos dois algoritmos é mais eficiente segundo essa medida. Para isso aplicou-se novamente o teste de Friedman com pós-teste de Nemenyi, dessa vez, com os melhores modelos dos dois algoritmos para cada medida. O resultado segue na Tabela~\ref{tab:melhoresmodelos}

%\begin{table}[!h]
%	\centering
%	
%	\begin{tabular}{|c|c|c|c|c|}
%
%		\hline
%		Medida & Algoritmo & \texttt{S} & \texttt{M} & \texttt{W}\\		
%		\hline
%		
%	
%		Acuracy		& C99 & 40 	& 11	& Sim \\ \hline
%		Precision	& C99 & 40 	& 11	& Sim \\ \hline
%		Pk			& C99 & 40 	& 11	& Sim \\ \hline
%		Windiff		& C99 & 40 	& 11	& Sim \\ \hline
%		F1			& C99 & 60 	& 9		& Sim \\ \hline
%		Recall		& C99 & 100 & 9		& Não \\ \hline
% 	
%	
%	\end{tabular}
%
%	\caption{Melhores modelos para cada medida segundo diagramas de diferença crítica.}
%	\label{tab:melhoresmodelos}	
%	
%\end{table}


%Na análise do diagrama de diferença crítica verificou-se que o algoritmo \textit{C99} apresenta melhor eficiência em todas as medidas e os valores das quatro primeiras os valores de \texttt{S}, \texttt{M} e \texttt{W} se repetiram, sugerindo uma configuração otimizada para o problema da segmentação de atas de reunião.




%tópicos retornando segmentos  que retorne segmentos coerentes 




%%%%%%%%%%
% Cálculo das medidas para cada modelo
%%%%%%%%%%

% --> Isso já foi falado no texto
% Pela comparação dos resultados com a segmentação fornecida pelos especialistas, calculou-se para cada modelo as medidas tradicionais acurácia, precisão, revocação, F-medida. 

% Além dessas, computou-se também as métricas mais aplicadas à segmentação textual P$_k$ e \textit{WindowDiff}.


















% ou seja limites não exatos podem ter um peso proporcional à proximidade, desde que estejam dentro da distancia K
%


% Baseando-se em medidas de avaliação tradicionais que utilizam a quantidade de acertos 


%	As medidas de avaliação tradicionalmente utilizadas na área de recuperação de informação como precisão e revocação trazem alguns problemas na avalização de segmentadores automáticos. 
%	
%	se considerarmos 

%	As medidas de avaliação tradicionalmente utilizadas na área de recuperação de informação como  
%
%Precisão é a proporção de limites identificados pelo algoritmo que são verdadeiros
%
%Revocação e a proporção limites verdadeiros que foram identificados pelo algoritmo

%corretamente identificados pelo algorítmo	
	 



%Para quantificar a eficiência dos algoritmos, segue uma revisão das principais métricas aplicáveis.



% para remoção de ruídos e ajustes 
